\section{Computational Cost and Accuracy-Reliability Trade-off}
\label{sec:cost_tradeoff}

The introduction of graph-based message passing improves forecasting accuracy,
but also increases model complexity and computational cost.
Table~\ref{tab:computational_cost} compares the Temporal model and the selected
Graph-Temporal model in terms of number of trainable parameters and training
time.

\begin{table}[t]
    \centering
    \caption{Computational cost of the Temporal and Graph-Temporal models.}
    \label{tab:computational_cost}
    \begin{tabular}{lrr}
        \toprule
        Model & Trainable Parameters & Training Time (min) \\
        \midrule
        Temporal GRU   & 40164 & 14.3 \\
        Graph-Temporal & 64868 & 34.0 \\
        \bottomrule
    \end{tabular}
\end{table}

As shown in Table~\ref{tab:computational_cost}, the Graph-Temporal model
contains 64868 trainable parameters compared with 40164 for the Temporal
model, corresponding to an increase of approximately 61.5\%. Training time also
increases from approximately 14.3 to 34.0 minutes, making the Graph-Temporal
model about 2.4 times more expensive to train in the considered experimental
setting. This additional cost results from the diffusion operations and from
the stacked graph-temporal blocks used to combine spatial and temporal
information.

The additional computational cost is accompanied by improved forecasting
accuracy. As shown in Section~\ref{sec:in_distribution_results}, the
Graph-Temporal model consistently outperforms the Temporal model in terms of
MAE, RMSE, and Pinball Loss, with the improvement in point forecasting
becoming more pronounced at longer horizons. It also produces sharper
prediction intervals, although with slightly lower empirical coverage.
Therefore, the accuracy gains provided by graph message passing come at the
cost of both increased computation and mild undercoverage.

Moreover, the stress-test analysis shows that this reliability gap becomes
more important under high congestion. Although the Graph-Temporal model
responds to difficult traffic conditions by producing wider intervals,
coverage decreases to 75.8\% at the 60-minute horizon. The model therefore
does not fully translate its increased uncertainty into reliable coverage
under distribution shift.

The CQR mitigation partially addresses this trade-off without increasing
model complexity or requiring retraining. Post-hoc calibration brings overall
coverage close to the nominal 80\% level while leaving the median point
forecast unchanged. Nevertheless, the remaining undercoverage under
high-congestion conditions shows that improved average calibration does not
necessarily imply equally reliable uncertainty estimates across different
traffic regimes.

Overall, the Graph-Temporal architecture provides a clear accuracy advantage
over the purely temporal baseline at the cost of greater computational
complexity. Its probabilistic forecasts are generally sharper and more
accurate in terms of quantile loss, but require additional calibration to
approach the desired coverage. The results therefore highlight that model
selection should consider not only point forecasting accuracy, but also
computational cost and the reliability of predictive uncertainty,
particularly under distribution shift.

